\section{把数据升到高维为什么有用}
\label{sec:13-Machine-Learning/09-Kernel-Methods-and-Gaussian-Processes/01}

产线上装了两只传感器，你要根据它们的读数判断设备是否异常。逻辑回归调了一周，验证集准确率始终在 \(50\%\) 上下打转，像是在抛硬币。翻开原始记录才看明白规律：两只读数同时偏高、或者同时偏低，设备都正常；一高一低才是异常。任何``给每只传感器配一个权重再相加''的模型都写不出这条规则。问题不在优化器，也不在数据量，在于假设类里根本不存在那个函数。

这一节从这类死局出发。先证明它确实是死局——不是没找到解，是解不存在；再给出补救：把数据映射到更高维，同一批点就能被一张平面分开。升维不是把问题变简单，它把原空间中的非线性关系改写成新特征上的线性关系，表达力和记住噪声的能力同步上升。所以``训练数据可分''只说明表示够丰富，不说明边界会泛化。

\subsection{一维就足以让线性模型失效}

先看一个连二维都不需要的例子。四个样本 \(x=-2,-1,1,2\)，标签按``绝对值是否大于 \(1.5\)''给出：\(-2\) 与 \(2\) 为正类，\(-1\) 与 \(1\) 为负类。一维线性分类器形如 \(f(x)=\sign(wx+b)\)，正类区域只能是半条直线：要么 \(x\) 大于某阈值，要么小于某阈值。而正类集合 \(\set{-2,2}\) 分居负类两侧，任何半直线只要同时含住 \(-2\) 与 \(2\)，就必然把 \(-1\) 或 \(1\) 一并吞进去。不存在任何 \((w,b)\) 能分对。

补救很廉价：加一个特征，令 \(\phi(x)=(x,x^2)\)。在新空间取 \(f(x)=x^2-\tfrac94\)，则 \(\abs{x}>1.5\) 时 \(f>0\)，否则 \(f<0\)，四个点全部分对。正类区域在原空间是``两端''，在新空间是一条水平线上方——非线性边界被折进了新特征里。

\subsection{四个点难住整张平面}

一维的失败也许还显得特殊。传感器那个例子是二维的 XOR，它是彻底的死局。四个样本的特征取值为 \(0\) 或 \(1\)，标签由 XOR 规则给出：

\[
(0,0)\mapsto -1,\quad
(1,1)\mapsto -1,\quad
(0,1)\mapsto +1,\quad
(1,0)\mapsto +1.
\]

任何线性分类器都形如 \(f(x)=\sign(w_1x_1+w_2x_2+b)\)。假设它分对了全部四个点，正类两样本给出

\[
w_2+b>0,\qquad w_1+b>0,
\]

负类两样本给出

\[
b<0,\qquad w_1+w_2+b<0.
\]

前两个不等式相加得到 \(w_1+w_2+2b>0\)，后两个相加得到 \(w_1+w_2+2b<0\)。同一个量不可能既正又负，因此不存在任何 \((w_1,w_2,b)\) 完成分类。这比``还没找到''强得多：两条对角线各属一类，任何直线都会切错至少一块。

失败的原因比失败本身更值得记住。XOR 的标签取决于两个特征的联合取值：单看 \(x_1\) 或 \(x_2\)，与标签毫无关系——在每个特征的任一取值下，正负类各占一半；合在一起却完全决定标签。线性模型只能表达``每个特征各记一笔账、加起来''的函数，而这里恰恰没有可以单独记的账。

\subsection{补上联合取值，平面重新够用}

既然缺的是``联合取值''，就把它补成一个新特征。定义

\[
\phi(x_1,x_2)=(x_1,\;x_2,\;x_1x_2).
\]

四个样本在新空间 \(\R^3\) 中成为

\[
\text{负类：}(0,0,0),\ (1,1,1);
\qquad
\text{正类：}(0,1,0),\ (1,0,0).
\]

取

\[
f(x)=1.5\,x_1+1.5\,x_2-3\,x_1x_2-1,
\]

逐点代入：\((0,0)\) 处为 \(-1\)，\((1,1)\) 处为 \(1.5+1.5-3-1=-1\)，\((0,1)\) 与 \((1,0)\) 处均为 \(1.5-1=0.5\)。负类全负、正类全正：在 \(\phi\) 给出的三维空间里，一张平面完成了分类。

下面这张图把两个空间并排放在一起。左边是原空间，分离面是一条双曲线；右边把坐标换成 \(u=x_1+x_2\) 与 \(v=x_1x_2\)，同一个 \(f\) 变成一条直线。

\begin{center}
\begin{tikzpicture}[x=1.55cm,y=1.55cm,>=stealth]

\begin{scope}
  \draw[->] (-0.35,0) -- (1.55,0) node[right] {\(x_1\)};
  \draw[->] (0,-0.35) -- (0,1.55) node[above] {\(x_2\)};
  \begin{scope}
    \clip (-0.32,-0.32) rectangle (1.5,1.5);
    \draw[thick,black!55] plot[domain=-0.25:0.40,samples=40]
      ({\x},{(1-1.5*\x)/(1.5-3*\x)});
    \draw[thick,black!55] plot[domain=0.62:1.45,samples=40]
      ({\x},{(1-1.5*\x)/(1.5-3*\x)});
  \end{scope}
  \fill (0,1) circle (2.6pt);
  \fill (1,0) circle (2.6pt);
  \draw[thick] (0,0) circle (2.6pt);
  \draw[thick] (1,1) circle (2.6pt);
  \node[font=\footnotesize] at (0.75,1.45) {原空间：边界是双曲线};
\end{scope}

\begin{scope}[shift={(4.6cm,0cm)}]
  \draw[->] (-0.35,0) -- (2.45,0) node[right] {\(u=x_1+x_2\)};
  \draw[->] (0,-0.55) -- (0,1.35) node[above] {\(v=x_1x_2\)};
  \draw[thick,black!55] (-0.2,-0.433) -- (2.4,0.867);
  \fill (1,0) circle (2.6pt);
  \draw[thick] (0,0) circle (2.6pt);
  \draw[thick] (2,1) circle (2.6pt);
  \node[font=\footnotesize,anchor=west] at (1.08,-0.22) {两个正类点重合};
  \node[font=\footnotesize] at (1.2,1.25) {升维后：边界是直线};
\end{scope}

\end{tikzpicture}
\end{center}

\emph{同一批数据、同一个判别函数，换个坐标就从弯曲的边界变成一条直线。}

图里发生了两件事。其一，``非线性问题''只是换到了另一个空间，并没有消失：右图那条直线拉回原空间就是左图的双曲线 \(1.5x_1+1.5x_2-3x_1x_2-1=0\)。其二，模型对参数仍然是线性的——写成 \(f=\ip{w}{\phi(x)}+b\)，\(w=(1.5,1.5,-3)\) 以线性方式出现。\hyperref[sec:13-Machine-Learning/02-Linear-Models/01]{``线性回归从投影到概率模型''}末尾强调过这一点：线性指对参数线性，对原始变量可以任意弯曲。升维没有改变模型的种类，只改变了模型看到的数据表示。

\subsection{维数越高，可分的标签方式越多}

XOR 不是孤立的巧例。Cover 定理（1965）给出一般性陈述：对处于一般位置的点，线性分类器能实现的二分类方式数量随特征维数上升；当特征维数接近样本数时，几乎所有标签方式都可分。这里的``一般位置''排除了特殊的仿射退化，计数结论也不是说随便一个非线性映射都会让数据可分。

直觉并不神秘。\(d\) 维空间中，\(n\) 个一般位置的点能被超平面实现的标签划分数目为

\[
2\sum_{k=0}^{d}\binom{n-1}{k},
\]

而全部标签方式共 \(2^n\) 种。\(d\) 远小于 \(n\) 时，可实现的比例微乎其微；\(d\) 接近 \(n-1\) 时，求和逼近 \(2^{n-1}\)，几乎每种标签划分都有某张超平面能实现。维数每增加一，空间就多出一点余地让点错开。

这句直觉必须立刻配上警告：分得开不等于分得好。若任意标签都能被分开，那么把噪声当标签也照样分得开——表达力的增长同时是记忆能力的增长。``可分性变容易''与``泛化变容易''是两件事，后者需要容量控制来平衡，\hyperref[ch:118]{``统计学习理论：训练误差对真风险有多乐观''}一章处理的正是这个差额。升维只解决了前半件事。

\subsection{显式构造特征的账单付不起}

既然升维有用，把``所有可能的交叉特征''都造出来似乎是一条路。设原始输入 \(x\in\R^d\)，构造全部不超过二次的特征：

\[
\phi(x)=\bigl(1,\;x_1,\dots,x_d,\;x_1^2,\;x_1x_2,\;\dots,\;x_d^2\bigr).
\]

特征总数为

\[
\binom{d+2}{2}=\frac{(d+2)(d+1)}{2}\approx \frac{d^2}{2}.
\]

\(d=1000\)（一张小图片、一段稀疏文本都远不止此）时约五十万维。这还只是二次：全部不超过五次的特征数为

\[
\binom{d+5}{5}\approx \frac{d^5}{120},
\]

\(d=1000\) 时约 \(8\times 10^{12}\) 维。一条样本的显式特征就要占据几十 TB 内存，内积、求解、存储全部不可行。手工挑少量交叉项可以缓解，却又回到``靠人猜哪些交互重要''的老问题——猜错的部分，模型永远补不回来。

账单的根源在于我们坚持显式构造 \(\phi(x)\)：先把每个样本变成一根高维向量，再在向量上做运算。如果后续算法根本不需要 \(\phi(x)\) 本身，只需要两个样本之间的某种相似程度，也许就能绕开整张账单。

\subsection{只算内积就够了}

回看已经学过的推导，会发现一个共同形态。\hyperref[sec:09-Convex-Optimization/07-Applications/03]{``最大间隔与支持向量机''}的对偶问题里，数据只以 \(x_i^{\trans}x_j\) 的形式出现；\hyperref[sec:13-Machine-Learning/02-Linear-Models/02]{``正则化从 Ridge 到 Lasso''}的一阶条件里，数据只以 \(X^{\trans}X\) 与 \(X^{\trans}y\) 的形式出现——同样全是内积的组合。把 \(x\) 换成 \(\phi(x)\) 后，这些算法只需要

\[
k(x,z)=\ip{\phi(x)}{\phi(z)},
\]

无须显式构造 \(\phi(x)\)。若有办法不算 \(\phi\) 而直接算出 \(k(x,z)\)，整张维度账单就换成一次便宜的函数求值：多项式特征的内积只要 \(O(d)\) 次运算，与特征空间是五十万维还是 \(8\times10^{12}\) 维无关。上图右侧的那两个新坐标之所以能被省掉，靠的正是这一步；下一节把它写成正式规则，并说明什么样的函数才有资格充当 \(k\)。
